\section{Parallel Treatments in the Works of St.~Thomas Aquinas}
\label{app:aquinas-parallels}

Aquinas revisited the angels across his career---from the early Parisian \emph{De ente et essentia} and \emph{De veritate} (c.~1256--1259), through the mature synthesis of the \emph{Summa theologiae} Prima Pars (c.~1266--1268) and Tertia Pars (c.~1272--1273), to his most extensive demonological treatise in \emph{De malo} q.~16 (c.~1270--1272). The table below maps the directly inspected parallel editions, highlighting every point where Aquinas sharpens, expands, or shifts an argument. Additional Aquinas treatises (\emph{Summa contra gentiles}, \emph{De spiritualibus creaturis}, \emph{De substantiis separatis}, and \emph{Scriptum super Sententiis} Book~II), not consulted here in full electronic editions, are excluded from verbatim claims and noted in Appendix~\ref{app:scope}.

{\footnotesize
\setlength{\tabcolsep}{2.5pt}
\renewcommand{\arraystretch}{1.15}
\begin{longtable}{>{\bfseries\raggedright\arraybackslash}p{0.18\linewidth}>{\raggedright\arraybackslash}p{0.14\linewidth}>{\raggedright\arraybackslash}p{0.24\linewidth}>{\raggedright\arraybackslash}p{0.38\linewidth}}
\toprule
\textbf{Work and Locus} & \textbf{\emph{ST} Parallel} & \textbf{Subject} & \textbf{Clarification, Development, or Shift} \\
\midrule
\endfirsthead
\toprule
\textbf{Work and Locus} & \textbf{\emph{ST} Parallel} & \textbf{Subject} & \textbf{Clarification, Development, or Shift} \\
\midrule
\endhead
\bottomrule
\endfoot
\emph{De ente et essentia}, chs.~4--5 & \ST{I, q.~50, aa.~2, 4} & Subsisting form, essence and \emph{esse}, individuation & Establishes that separate substances lack matter yet remain composed of form (essence, potency) and received \emph{esse} (act); refutes universal hylomorphism and derives one angel per species. \\
\addlinespace[2pt]
\emph{De veritate}, q.~8, aa.~1--4 & \ST{I, q.~12; q.~56, a.~3; q.~57, a.~5} & Beatific vision, \emph{lumen gloriae}, natural knowledge of God & Declares denial of immediate vision of God's essence heretical (a.~1); in a.~3~ad~17 treats natural angelic knowledge of God as intermediate via impressed species, whereas \ST{I, q.~56, a.~3} aligns it more closely with specular mirror-knowledge through the angel's own essence. \\
\addlinespace[2pt]
\emph{De veritate}, q.~8, aa.~5--15 & \ST{I, qq.~55--58} & Media of knowledge, singulars, futures, hearts, simultaneity & Develops the architectonic analogy for universal innate species (a.~10), surveys four opinions on singulars (a.~11), and expounds Dionysius's circular, straight, and oblique intellectual motions (a.~15~ad~3). \\
\addlinespace[2pt]
\emph{De veritate}, q.~8, aa.~16--17 & \ST{I, q.~58, aa.~6--7} & Morning and evening knowledge & Strictly separates morning and evening knowledge by uncreated vs.\ created cognitive medium (a.~16), whereas \ST{I, q.~58, a.~7} accommodates unity when things in proper nature are referred to the Word. \\
\addlinespace[2pt]
\emph{De veritate}, q.~9, aa.~1--7 & \ST{I, qq.~106--107} & Angelic illumination, purgation, and speech (\emph{locutio}) & Explicitly states illumination infuses no new species (a.~1~ad~10); distinguishes \emph{locutio} from \emph{illuminatio} by object presence vs.\ intellectual strength (a.~5), whereas \ST{I, q.~107, a.~2} sharpens the line metaphysically by causal root (First Truth vs.\ created will). \\
\addlinespace[2pt]
\emph{De malo}, q.~16, a.~1 & \ST{I, q.~50, a.~1; q.~51, a.~1} & Whether demons have natural aerial bodies & Opens with Augustine's candor that aerial bodies do not strictly affect Christian faith; allows in ad~3 and ad~7 that Dionysius (``in many things a follower of Plato'') and Damascene may have entertained subtle bodies---a freer historical concession than \ST{I, q.~51, a.~1}. \\
\addlinespace[2pt]
\emph{De malo}, q.~16, aa.~2--3 & \ST{I, q.~63, aa.~1--4} & Evil by will, non-consideration, coveting equality & Clarifies that the angel sinned through non-consideration of the higher divine rule without prior speculative error (a.~2), coveting supernatural beatitude by natural power alone rather than by grace (a.~3). \\
\addlinespace[2pt]
\emph{De malo}, q.~16, a.~4 & \ST{I, q.~63, aa.~5--6} & Sin in the first instant of creation & \textbf{Major doctrinal shift:} \ST{I, q.~63, a.~5} argues from efficient causality (allowing instant-1 merit, ad~3--4); \emph{De malo} argues from non-simultaneous cognitive actuation of natural vs.\ supernatural objects, concluding that in instant~1 the angel was neither blessed nor sinful, moving solely toward God as Author of nature (ad~14--15). Notes in ad~16 that humanity would have been created even had no angel fallen. \\
\addlinespace[2pt]
\emph{De malo}, q.~16, aa.~5--6 & \ST{I, q.~64, aa.~1--2} & Irrevocable obstinacy and darkened intellect & Gives a threefold analysis of \emph{liberum arbitrium}; notes in a.~5~ad~13 that God could convert a demon by \emph{potentia absoluta} though incongruous with angelic nature; distinguishes simple nescience from proud speculative and practical error (a.~6). \\
\addlinespace[2pt]
\emph{De malo}, q.~16, aa.~7--8 & \ST{I, q.~57, aa.~3--4} & Demonic knowledge of futures and heart-thoughts & Supplies the watchtower-of-eternity image (a.~7) and distinguishes stored intelligible species (accessible to demons) from the voluntary \emph{usus speciei} (hidden from all creatures, a.~8). \\
\addlinespace[2pt]
\emph{De malo}, q.~16, aa.~9--12 & \ST{I, qq.~110--111, 114} & Transmutation, local motion, senses, intellect & Clarifies that demons cannot move celestial spheres or the whole earth (a.~10), act on imagination by moving bodily spirits/humors (a.~11), and can cloud human reason via phantasms only as permitted by God (a.~12). \\
\addlinespace[2pt]
\emph{ST} III, q.~8, a.~4 & \ST{I, q.~108, aa.~1, 8; q.~113, a.~4} & Christ as Head of the angels & Establishes Christ's headship even as man through generic conformity of soul with angelic spirit, heavenly ecclesial unity of comprehensors, and instrumental efficacy of Christ's sacred humanity (expanding \emph{Super Col.}~2:10). \\
\addlinespace[2pt]
\emph{ST} III, q.~30, aa.~1--4 & \ST{I, q.~51, a.~2; q.~57, a.~5; q.~112, a.~2} & The Annunciation to Mary & Synthesizes Gabriel's rank (highest of archangels, reconciling Gregory and Dionysius), Mary's consent in lieu of human nature, waking bodily vision, and Antony's criterion of peace succeeding fear. \\
\addlinespace[2pt]
\emph{Super Eph.} \& \emph{Super Col.} (Dessain 1857 OCR finding aid) & \ST{I, q.~57, a.~5; q.~108, a.~6; III, q.~8, a.~4} & Pauline angelic hierarchies and Eph~3:10 & Cross-references Eph~1:21 and Col~1:16 on Dionysius vs.\ Gregory; resolves Eph~3:10 (Church makes wisdom known \emph{in} angels rather than teaching them from below). \\
\end{longtable}
}
